\section{Propuesta de solución}

Reuniendo todo lo mencionado ante la problemática presente en la sociedad colombiana, que de paso no será algo que se solucione en el corto plazo y el país necesita una solución al menos temporal mientras todo retoma su ``curso normal''.

La solución planteada es implementar un modelo IA que tenga un aprendizaje automatizado supervisado (CNN) y no supervisado (KNN), el cual se entrenará previamente con la mayor cantidad de datos que logremos conseguir, en donde el paciente sube la foto de un medicamento, bajo un contexto que fue un medicamento que utilizó previamente más ahora que lo volvió a buscar se dio cuenta que ya se le acabó, entonces, el modelo se encargará de darle al paciente/persona información sobre ese medicamento (su tipo, si es analgésico, antibiótico o antiinflamatorio, y de paso información como tal del medicamento). De igual manera, se le dará al paciente recomendaciones del fármaco que colocó, sea que tenga el mismo componente a que tenga un componente parecido, pero que, de todas formas, le funcione a la persona para que lo pueda conseguir en caso de que el que tenía que se le agotó, se le complique conseguirlo.

\subsection{Arquitectura del sistema}

Mediante el conocimiento que tenemos sobre implementar modelos IA en código, inicialmente vamos a empezar a implementar todo con las librerías Scikit-learn + TensorFlow/PyTorch en Python, para no tener que definir todo desde cero. Igualmente, en caso de que sea necesario implementarlo desde cero, se hará definiendo todas las funciones básicas para que el modelo funcione de manera correcta.

Tener en cuenta que toca tener metadata asociada, pero toda la búsqueda de similitud es trabajo 100\% de ML, en donde:

\subsubsection{Durante entrenamiento}

\begin{itemize}
	\item Guardas los embeddings de todos los medicamentos del training set
	\item Guardas metadata asociada (nombre, componente, usos)
\end{itemize}

\subsubsection{Durante inferencia}

\begin{itemize}
	\item CNN extrae embedding del medicamento nuevo
	\item KNN busca en el espacio de embeddings (no en una BD tradicional)
	\item KNN retorna los K vecinos + su metadata
\end{itemize}

La clave está en que el KNN trabaja en el espacio latente aprendido por la CNN, no en una base de datos relacional.

\pagebreak